\chapter{Identifiers for digital objects}
\label{ch:identifiers}

The previous chapter left us with a quietly remarkable fact: in the Software Heritage data model,
every artifact already \emph{has} a name, computed from its own content. Before we put that name
to work, we should be clear about what a name is \emph{for} --- because almost everything people
believe about digital identifiers is tangled in a confusion that a few minutes of care can
dissolve.

\takeaway{Identification is one thing; location, metadata, organization, and curation are four
others. Keep them apart.}

\section{Five concerns hiding behind one word}

When we say that a digital object \enquote{has a DOI}, we are quietly running five different
concerns together. Pulling them apart is the single most clarifying move in this whole subject:

\begin{description}[leftmargin=!,labelwidth=2.5cm,font=\bfseries\sffamily]
  \item[Identification] giving the object a unique name --- the identifier proper.
  \item[Location] a way to \emph{access} or resolve the object. The ISBN names a book; the library
        catalogue tells you which shelf it sits on. Two different jobs.
  \item[Metadata] what is \emph{known about} the object --- authors, title, date --- usually held
        in a registry.
  \item[Organization] who keeps the system coherent: the governance, the registry operator, the
        institutions that keep the lights on.
  \item[Curation] the \emph{quality} of the metadata --- how it is checked and maintained, and who
        decides which objects deserve metadata, or even a name, at all.
\end{description}

Only the first is \emph{identification}. The other four are real and important --- but they are
not the same thing, and most arguments about identifiers turn out to be arguments about one of the
latter four, wearing identification's clothes. One example makes this unforgettable.

\section{``Crossref DOIs are better than DataCite DOIs''}

You will hear this said, and it is worth taking seriously, because it is at once obviously true
and strictly nonsensical --- depending on which of our five concerns you have in mind.

Strictly, it is nonsense. A DOI is a DOI. Crossref and DataCite are both \emph{registration
agencies} accredited by the same International DOI Foundation; both mint identifiers resolved by
the same underlying Handle System~\autocite{Arms01,RFC3650}, through the same \texttt{doi.org}.
They are so much one system that when \texttt{doi.org} suffered an outage in 2015, it took down
\emph{every} agency's DOIs at once. Hand someone a bare DOI and, as a rule, they cannot even tell
you which agency registered it. At the level of \emph{identification}, then, \enquote{better} has
no meaning: the identifiers are the same kind of object, named and resolved in exactly the same
way~\autocite{doifoundation_ra}.

And yet the sentence is not foolish; the people who say it mean something real. They are talking
about \emph{organization} and \emph{curation}. Crossref grew up around the scholarly literature,
DataCite around research data and software; the two run different metadata schemas, serve
different communities, and accumulate metadata of differing richness and quality. But that
richness, as the two organizations themselves put it, \enquote{hinges on what depositing members
actively provide, rather than being automatically generated by the DOI assignment process
itself}~\autocite{crossrefdatacite_metadata}. The quality lives in the organization and its
curation, not in the identifier. So the famous sentence is true and false at the same time: false
about identification, true about everything else --- and it sounds paradoxical only because all
five concerns were crammed into one word. Pull them apart, and the paradox evaporates.

\section{Intrinsic versus extrinsic identifiers}

There is a deeper lesson hiding in that example, and Norman Paskin --- long the voice of the DOI
itself --- put his finger on it: \enquote{Digital Object Identifier} is best read as \emph{digital
identifier of an object}, not \emph{identifier of a digital object}~\autocite{paskin2010digital}.
That distinction is the whole game.

An \emph{extrinsic} identifier --- a DOI, a Handle, an ARK --- is a digital identifier \emph{of}
something, assigned by an organization and recorded in a register. It says nothing about the bytes
it names; it is a pointer whose meaning is maintained, on goodwill, by several parties at once. The
resolution can be repointed; the content at the far end can change; and there is \emph{no intrinsic
way to notice}~\autocite{swhipres2018}. Anyone who has clicked a DOI and met \enquote{not found}
has felt the consequence. Such an identifier delegates to an organization not only \emph{location}
but \emph{identification} itself: without the register, it means nothing.

An \emph{intrinsic} identifier --- a Git commit hash, a \textsc{swhid} --- is an identifier
\emph{of} a digital object, computed from the object's own bytes. It needs no organization to
\emph{mean} something; it needs one only, and separately, to help you \emph{locate} what it names.
As we saw in Chapter~\ref{ch:foundations}, this is exactly what the Merkle construction buys:
identifiers that are gratis, that guarantee \emph{integrity}, and that need \emph{no
middle-man}~\autocite{cise-2020-doi}. Map this back onto our five concerns and the elegance is
plain: an extrinsic identifier entangles identification with location, metadata, organization, and
curation all at once; an intrinsic identifier \emph{separates} identification from all four. That
separation is no technicality. It is the difference between a name you must trust someone to keep,
and a name anyone can check.

\takeaway{An intrinsic identifier is computed from the content; an extrinsic one is assigned by a
registry.}

\section{The SWHID}

The intrinsic identifier for source code is the \textsc{swhid} --- the \emph{Software Hash
Identifier} --- and it is simply the name of a node in the Merkle DAG of the previous chapter. In
its core form it reads \swhid{swh:1:\textit{type}:\textit{hash}}, where the type is one of
\texttt{cnt} (a file's content), \texttt{dir} (a directory), \texttt{rev} (a revision, i.e.\ a
commit), \texttt{rel} (a release), or \texttt{snp} (a snapshot of a whole repository). Thus
\swhid{swh:1:cnt:94a9ed02\ldots} \emph{is} the text of the GPL\,v3 licence, wherever it appears,
and \swhid{swh:1:rel:22ece559\ldots} is release 2.3.0 of the Darktable project --- not a record
pointing \emph{at} them, but a name computed \emph{from} them, that you can recompute yourself.

Identification, though, is not location, and the \textsc{swhid} keeps the two cleanly apart. When
you do want to point a reader into the archive --- to say not merely \emph{which} code but
\emph{where it came from, and in what context} --- you add qualifiers: \texttt{origin} (the URL the
code was harvested from), \texttt{visit}, \texttt{anchor}, \texttt{path}, and \texttt{lines}. The
core name still identifies; the qualifiers merely locate. Appendix~\ref{app:swhid-syntax} gives the
full grammar.

None of this is a private convention. The \textsc{swhid} is an international standard, ISO/IEC
18670:2025, recognised by SPDX, registered with IANA as the \texttt{swh:} URI scheme, and recorded
as Wikidata property P6138~\autocite{isoiec18670}. And it is not theoretical for this very note: as
the Colophon explains, this document is named and verified by its own \textsc{swhid} --- which is
also exactly why it cannot contain that identifier inside itself.

\begin{handson}[Hands-on: watch the name follow the bytes]
\begin{enumerate}[nosep,leftmargin=*]
  \item Compute a content \textsc{swhid} for any file: \texttt{swh identify file.c}.
  \item Change a single byte and recompute --- the identifier changes completely.
  \item Restore the byte; the original identifier returns, to the character. No registry was ever
        consulted, and no authority was asked --- the name simply follows the bytes.
\end{enumerate}
\end{handson}

We now have what the rest of this note needs: a name we can trust, that identifies code exactly,
that anyone can verify without asking permission, and that --- when we choose --- can also tell us
where the code lived. The remaining chapters are about putting that name to work: archiving the
code, referencing it at the right granularity, describing it, and citing it. That practice is
\textsc{ardc}, and it is next.
